(1) 将 \(P=RT/(v-b)-a/v^2\) 代入 \(h=u+Pv\)，得到 \(\boxed{h=C_vT+RTv/(v-b)-2a/v}\)。绝热多孔塞节流在忽略动能改变时保持焓，而自由膨胀隔热且无外功时保持内能，不能把两种过程混作等温。

(2) 利用状态方程在恒压条件隐式求导 \(v_T|_P=[R/(v-b)]/[RT/(v-b)^2-2a/v^3]\)。对 \(C_P>0\) 的稳定区域，\(\mu_{JT}\) 的符号由 \(Tv_T-v\) 决定；反转线要求 \(T v_T=v\)，代入方程即可消去偏导。

(3) 反转条件整理为 \(TRb/(v-b)^2=2a/v^2\)，即 \(\boxed{T_{\rm inv}(v)=\frac{2a}{Rb}(1-b/v)^2}\)，对应压强 \(\boxed{P_{\rm inv}(v)=a(2v/b-3)/v^2}\)。在 \(v\to\infty\) 的低压端，\(T_{\rm inv}\to2a/(Rb)\)、\(P_{\rm inv}\to0\)；只有对应物理单相稳定区才可实验观测。

(4) 对 \(v\gg b\) 作Virial展开，可得 \(\mu_{JT}\simeq[2a/(RT)-b]/C_P\)。吸引作用 \(a\) 在温度较低时造成节流降温，排斥排除体积 \(b\) 则贡献相反符号。真实气体的 \(C_P\)、多体相互作用和相变区域会修正简化反转线。
